\begin{code}[hide]
module slides.Tsinghua-PL-seminar.sections.newqupit where

open import Data.Nat using (ℕ)
open import Level using (0ℓ)
open import Algebra.Bundles using (Group)
open import Notations using (₁₊ ; ₂₊ ; ₃₊)
open import Word.Base using (Word ; WRel ; ε ; _•_)
open import Presentation.Definitions using (_IsPresentationOf_)

infix 4 _SRel,_===_ _QRel,_===_
infixr 8 _^_
infix 9 S^_ Z^_ CZ^_ XMg^_ g^_
infixl 10 _↑ _↓
infixl 7 _*_
infixl 6 _+_
infix 8 -_

postulate
  Gen : ℕ → Set
  p p-1 g k 1/2 : ℕ
  S H Z CZ Ex CX CZ02 XMg XM₋₁ : ∀ {n} → Word (Gen n)
  XM : ℕ → ∀ {n} → Word (Gen n)
  _^_ : ∀ {n} → Word (Gen n) → ℕ → Word (Gen n)
  S^_ Z^_ CZ^_ XMg^_ : ∀ {n} → ℕ → Word (Gen n)
  g^_ -_ : ℕ → ℕ
  _*_ _+_ : ℕ → ℕ → ℕ
  _↑ _↓ : ∀ {n} → Word (Gen n) → Word (Gen n)
  _QRel,_===_ : (n : ℕ) → WRel (Gen n)
  Sp-group Pauli⋊Sp : ℕ → Group 0ℓ 0ℓ

private variable
  n : ℕ
\end{code}
%% Part 6b — qupit Clifford circuits (branches main, qupit, ...).

\begin{frame}{Qupit Clifford circuits in all odd prime dimensions \New}
\begin{columns}[c]
\begin{column}{0.40\textwidth}
\small
\textbf{Bian, Li, Ross, van de Wetering, Zhao}, \emph{A Complete and Natural
Rule Set for Multi-Qudit Clifford Circuits in All Odd Prime Dimensions}
(QPL'26). Its verification is a separate paper in preparation.
\begin{itemize}
  \item gates $H$, $S$ (one qupit), $CZ$ (two), any width $n$, any odd prime $p$;
  \item read \alert{modulo scalars}, the circuits present
        \[\text{Pauli}(n)\rtimes Sp(2n,\mathbb{Z}_p);\]
  \item rules over $p$, a \emph{primitive root} $g$, and $\frac12\in\mathbb{Z}_p$;
  \item 16 rules in the paper (Figure 1); 15 here: $(SH)^3=\varepsilon$
        becomes a theorem.
\end{itemize}
\end{column}
\begin{column}{0.58\textwidth}
\renewcommand{\AgdaCodeStyle}{\scriptsize}
\begin{code}
data _SRel,_===_ : (n : ℕ) → WRel (Gen n) where
  order-S :       (₁₊ n) SRel,  S ^ p === ε
  order-H :       (₁₊ n) SRel,  H ^ 2 === XM₋₁
  M-power : ∀ k → (₁₊ n) SRel,  XMg^ k === XM (g^ k)
  semi-MR :       (₁₊ n) SRel,  XMg • S^ (g * g) • Z^ ((g * g + - g) * 1/2)
                                === S • XMg
  comm-HHSHHS :   (₁₊ n) SRel,  H • H • S • H • H • S === S • H • H • S • H • H
  order-CZ :      (₂₊ n) SRel,  CZ ^ p === ε
  order-Ex :      (₂₊ n) SRel,  Ex ^ 2 === ε
  comm-CZ-S↑ :    (₂₊ n) SRel,  CZ • S ↑ === S ↑ • CZ
  semi-M↑CZ :     (₂₊ n) SRel,  XMg ↑ • CZ^ g === CZ • XMg ↑
  semi-Ex-S↑ :    (₂₊ n) SRel,  Ex • S ↑ === S ↓ • Ex
  semi-Ex-H↑ :    (₂₊ n) SRel,  Ex • H ↑ === H ↓ • Ex
  blake-c12 :     (₂₊ n) SRel,  (S ^ p-1) ↑ • (S ^ p-1) ↓ • CX ^ p-1 • S ↓ • CX === CZ
  yang-baxter :   (₃₊ n) SRel,  Ex ↑ • Ex ↓ • Ex ↑ === Ex ↓ • Ex ↑ • Ex ↓
  cz-slide :      (₃₊ n) SRel,  Ex ↓ • Ex ↑ • CZ === CZ ↑ • Ex ↓ • Ex ↑
  semi-CX↑-CZ↓ :  (₃₊ n) SRel,  CZ ↓ • CX ↑ === CZ02 • CX ↑ • CZ ↓
\end{code}
{\scriptsize \texttt{ProjectiveClifford/Qupit/Paper-V1/Syntactics.agda} (\Branch{qupit}), one comment elided}
\end{column}
\end{columns}
\end{frame}

\begin{frame}{The rules as circuits (one and two qupits) \New}
\begin{columns}[T]
\begin{column}{0.46\textwidth}\tiny
\begin{tabular}{@{}r@{\;\,}r@{\,}c@{\,}l@{}}
order-S & \qfig{pv0-order-S-l} & $\equiv$ & \qfig{pv0-order-S-r}\\[2pt]
order-H & \qfig{pv0-order-H-l} & $\equiv$ & \qfig{pv0-order-H-r}\\[2pt]
M-power & \qfig{pv0-M-power-l} & $\equiv$ & \qfig{pv0-M-power-r}\\[2pt]
semi-MR & \qfig{pv0-semi-MR-l} & $\equiv$ & \qfig{pv0-semi-MR-r}\\[2pt]
order-CZ & \qfig{pv0-order-CZ-l} & $\equiv$ & \qfig{pv0-order-CZ-r}\\[2pt]
order-Ex & \qfig{pv0-order-Ex-l} & $\equiv$ & \qfig{pv0-order-Ex-r}\\[2pt]
order-SH & \qfig{pv0-order-SH-l} & $\equiv$ & \qfig{pv0-order-SH-r}
\end{tabular}
\end{column}
\begin{column}{0.52\textwidth}\tiny
\begin{tabular}{@{}r@{\;\,}l@{}}
comm-CZ-S↑ & \qfig{pv0-comm-CZ-Sup}\\[2pt]
semi-M↑CZ & \qfig{pv0-semi-Mup-CZ}\\[2pt]
semi-Ex-S↑ & \qfig{pv0-semi-Ex-Sup}\\[2pt]
semi-Ex-H↑ & \qfig{pv0-semi-Ex-Hup}\\[2pt]
comm-HHSHHS & \qfig[0.5]{pv0-comm-HHSHHS}
\end{tabular}
\end{column}
\end{columns}
\vspace{2pt}
{\scriptsize Figure 1 of the QPL'26 paper (the multiplier spelled over $R=S\,Z^{1/2}$);
figures generated by Cir2Tikz from the Agda terms.}
\end{frame}

\begin{frame}{\ldots\ three qupits, and the symplectic group underneath \New}
\begin{columns}[c]
\begin{column}{0.50\textwidth}\tiny
\begin{tabular}{@{}r@{\;\,}l@{}}
blake-c12 & \qfig{pv0-blake-c12}\\[3pt]
yang-baxter & \qfig{pv0-yang-baxter}\\[3pt]
cz-slide & \qfig{pv0-cz-slide}\\[3pt]
semi-CX↑-CZ↓ & \qfig{pv0-semi-CXup-CZdn}
\end{tabular}
\end{column}
\begin{column}{0.48\textwidth}
\small $Sp(2n,\mathbb{Z}_p)$ \textbf{by a coset tower} (\texttt{Examples/Groups/Symplectic}, 35k lines):
\begin{itemize}
  \item one level per wire, Selinger-style staircases as digits;
  \item one wire: $\mathbb{Z}_p\times\mathbb{Z}_p^{*}\times(\varepsilon\mid HS^a)$, i.e.
        $p\,(p{-}1)(p{+}1) = |Sp(2,p)|$ normal forms;
  \item \alert{67 $\to$ 18}: normalization uses 67 relations; a small set of
        natural ones derives them all.
\end{itemize}
\begin{code}
-- MainTheorems, Symplectic-Simplified-Theorems
simplified-presentation :
  ∀ n → (n QRel,_===_) IsPresentationOf (Sp-group n)
\end{code}
\begin{code}[hide]
simplified-presentation = λ _ → _
\end{code}
\end{column}
\end{columns}
\end{frame}

\begin{frame}{Assembled, not re-proved \New}
\begin{columns}[c]
\begin{column}{0.52\textwidth}
\centering
\begin{tikzpicture}[node distance=5mm and 6mm,>={Stealth[length=1.8mm]},
  bx/.style={draw=accent,fill=soft,rounded corners=2pt,font=\scriptsize,align=center,minimum height=8mm,text width=26mm},
  nb/.style={draw=newC,fill=softnew,rounded corners=2pt,font=\scriptsize,align=center,minimum height=8mm,text width=26mm}]
  \node[bx] (pl) {ProjectivePauli\\{\tiny $(\mathbb{Z}_p\times\mathbb{Z}_p)^n$, the paper's Pauli}};
  \node[bx,right=of pl] (sp) {Symplectic.Simplified\\{\tiny coset tower, 18 rules}};
  \node[nb,below=9mm of $(pl)!0.5!(sp)$] (sd) {SemiDirect\\{\tiny $\text{Pauli}\rtimes Sp$, generic $\rtimes$ theorem}};
  \node[nb,below=of sd] (v1) {Simplified-V1\\{\tiny the Pauli calculus lives here}};
  \node[nb,below=of v1] (p0) {Paper-V0 (16 rules)\\{\tiny Figure 1}};
  \node[nb,below=of p0] (p1) {Paper-V1 (15 rules)};
  \draw[->,thick] (pl) -- (sd); \draw[->,thick] (sp) -- (sd);
  \draw[<->,thick,newC] (sd) -- node[right,font=\tiny]{iso} (v1);
  \draw[<->,thick,newC] (v1) -- node[right,font=\tiny]{iso, exports \aid{v1⇒pap}} (p0);
  \draw[<->,thick,newC] (p0) -- node[right,font=\tiny]{identity on words} (p1);
  \node[font=\tiny,gray,align=left,anchor=west] at ($(sd.east)+(1mm,0)$) {+ two well-definedness\\proofs of the action};
\end{tikzpicture}
\end{column}
\begin{column}{0.46\textwidth}
\small Every layer above the base is \alert{an isomorphism plus a composition};
no rule set re-proves what another has.
\begin{code}
-- MainTheorems, Qupit-Clifford-Theorems
clifford-presentation :
  ∀ n → (n QRel,_===_) IsPresentationOf (Pauli⋊Sp n)
\end{code}
\begin{code}[hide]
clifford-presentation = λ _ → _
\end{code}
\vspace{-2pt}
{\scriptsize (\aid{PapSyn.Clifford-Relations.\_QRel,\_===\_} in the source)}
\begin{itemize}
  \item reuses the paper's machinery verbatim: coset towers, the
        word-valued semidirect product, \aid{StarIsomorphism};
  \item derivability results: two Pauli-vs-$CZ$ axioms dropped (18 $\to$ 16);
        $(SH)^3=\varepsilon$ a theorem in V1;
  \item $\approx$22k lines for the four rule sets.
\end{itemize}
\end{column}
\end{columns}
\end{frame}

\begin{frame}{Getting the scalars back: exact Clifford groups \New}
\begin{columns}[c]
\begin{column}{0.50\textwidth}
\small The Clifford group is a \alert{central extension} of the projective one by
$\langle\omega\rangle$. Add $\omega$, $\omega^p=\varepsilon$, centrality, and
\alert{correct} each relation by the scalar it picks up --- only one needs it:
\vspace{4pt}

\centering\qfig[0.85]{sc-order-SH}

{\scriptsize ($\frac18$ the inverse of 8 in $\mathbb{Z}_p$)}
\vspace{4pt}

\raggedright
\small \texttt{Clifford/Qupit}: \aid{presentation-exact}, against the central
extension $A\times_\gamma H$ with $\gamma$ the Weyl cocycle $\frac12\,\text{sform}$;
new \texttt{ForStdlib}: \aid{CentralExtension}, \aid{FactorSetExtension}
(Rotman 9.8), \aid{CentralProduct}.
\end{column}
\begin{column}{0.47\textwidth}
\small \textbf{Qubits, after Selinger} (LMCS 2015), \texttt{ProjectiveClifford/Qubit},
\texttt{Clifford/Qubit}:
\begin{itemize}
  \item Figure 8 modulo scalars presents $V\!Sp(n)$ (symplectic map + $\mathbb{Z}_4$
        phase function); $\cong$ the Clifford action on signed Paulis;
  \item an \alert{extension} presentation $1\to\mathbb{Z}_8\to\text{Exact}(n)\to
        \text{CMS}(n)\to1$, isomorphic to Figure 8;
  \item that $\omega$ has order exactly 8 needs a model: $2^n\times2^n$
        matrices over $\mathbb{Z}/17$ with $i=4$, $\frac1{\sqrt2}=14$,
        $\omega=2$ --- every relation by \aid{refl} on tries.
\end{itemize}
\begin{newdevbox}
\small A first, small \alert{matrix} semantics --- next: real ones.
\end{newdevbox}
\end{column}
\end{columns}
\end{frame}
